\documentclass[fontsize=10pt,oneside,parskip=half]{scrreprt}
\usepackage{ulille-rapport}
\usepackage{amsmath}
\usepackage{float}
\usepackage{booktabs}
\usepackage{tabularx}
\usepackage{listings}

\newcommand{\ic}[1]{{\setlength{\fboxsep}{2pt}\colorbox{BeigeSerein}{\texttt{#1}}}}

\subject{SAE 2.01 - 2.02 : Rapport graphes}
\title{Seconde Version}
\author{Yann Renard \\ Tom Cox \\ Alexandre Sorel}
\direction{Équipe E5}

\lstdefinestyle{java}{
  language=Java,
  basicstyle=\ttfamily\small,
  keywordstyle=\bfseries\color{IndigoInclusif},
  commentstyle=\itshape\color{GrisStatutaire},
  stringstyle=\color{VertAvenir},
  backgroundcolor=\color{BeigeSerein},
  frame=single,
  framesep=6pt,
  rulecolor=\color{BleuHorizon},
  breaklines=true,
  breakatwhitespace=true,
  showstringspaces=false,
  tabsize=2,
  numbers=left,
  numberstyle=\tiny\color{GrisStatutaire},
  numbersep=8pt,
  xleftmargin=20pt,
  xrightmargin=4pt,
}

\begin{document}
\selectlanguage{french}

\maketitle
\pagecolor{Blanc}
{%
  \setlength{\parskip}{0pt}%
  \tableofcontents%
}

% ─────────────────────────────────────────────
\chapter{Version 1 : Un seul moyen de transport}
% ─────────────────────────────────────────────

% ──────────────────────────────
\section{Exemple}
% ──────────────────────────────

Pour illustrer la Version~1, nous avons choisi un réseau de transport reliant
cinq villes du nord de la France : Lille, Dunkerque, Cambrai, Paris et
Crépy-en-Valois. Ce réseau comprend dix-huit connexions réparties sur trois modes
de transport (train, bus et avion) avec des données de coûts estimées à partir
des vitesses moyennes de chaque mode et des barèmes d'émissions carbone de
l'ADEME.

\noindent\textbf{Tableau 1 : Données complètes du réseau de transport}

% ──────────────────────────────────────────────────────────────────────────────
% Tableau 1 : Données complètes du réseau de transport
% 
% Méthodologie de calcul :
% - TRAIN : vitesse moyenne 180 km/h, prix ~0.18 €/km, CO2 ~4 g/km/voyageur
% - BUS   : vitesse moyenne 90 km/h,  prix ~0.12 €/km, CO2 ~21 g/km/voyageur  
% - AVION : vitesse 800 km/h + 120 min temps d'enregistrement/aéroport
%          Prix fixe selon marché, CO2 ~85 g/km/voyageur (supposé 150 places)
% ──────────────────────────────────────────────────────────────────────────────

\begin{table}[H]
  \centering
  \begin{tabular}{lllrrr}
    \toprule
    \textbf{Départ} & \textbf{Arrivée} & \textbf{Mode}
      & \textbf{Prix (€)} & \textbf{CO\textsubscript{2} (kg)}
      & \textbf{Durée (min)} \\
    \midrule
    % Lille ↔ Dunkerque : 90 km
    Lille             & Dunkerque         & TRAIN  & 18  & 0{,}4  & 30  \\ % 90÷180×60 = 30 min
    Dunkerque         & Lille             & TRAIN  & 18  & 0{,}4  & 30  \\
    
    % Lille ↔ Cambrai : 70 km
    Lille             & Cambrai           & TRAIN  & 14  & 0{,}3  & 23  \\ % 70÷180×60 ≈ 23 min
    Cambrai           & Lille             & TRAIN  & 14  & 0{,}3  & 23  \\
    
    % Cambrai ↔ Paris : 160 km
    Cambrai           & Paris             & TRAIN  & 35  & 0{,}6  & 53  \\ % 160÷180×60 ≈ 53 min
    Paris             & Cambrai           & TRAIN  & 35  & 0{,}6  & 53  \\
    
    % Paris ↔ Crépy-en-Valois : 80 km
    Paris             & Crépy-en-Valois   & TRAIN  & 12  & 0{,}3  & 27  \\ % 80÷180×60 ≈ 27 min
    Crépy-en-Valois   & Paris             & TRAIN  & 12  & 0{,}3  & 27  \\
    
    % Lille ↔ Paris : 220 km
    Lille             & Paris             & TRAIN  & 55  & 0{,}9  & 73  \\ % 220÷180×60 ≈ 73 min
    Paris             & Lille             & TRAIN  & 55  & 0{,}9  & 73  \\
    
    % Dunkerque ↔ Cambrai : 100 km
    Dunkerque         & Cambrai           & TRAIN  & 22  & 0{,}4  & 33  \\ % 100÷180×60 ≈ 33 min
    Cambrai           & Dunkerque         & TRAIN  & 22  & 0{,}4  & 33  \\
    
    % Dunkerque → Crépy-en-Valois : 230 km (via Paris)
    Dunkerque         & Crépy-en-Valois   & TRAIN  & 45  & 0{,}9  & 77  \\ % 230÷180×60 ≈ 77 min
    
    % Routes en BUS (90 km/h + 5 min temps d'attente)
    Dunkerque         & Cambrai           & BUS    & 12  & 2{,}1  & 71  \\ % 100÷90×60 + 5 ≈ 71 min
    Cambrai           & Crépy-en-Valois   & BUS    & 18  & 2{,}5  & 85  \\ % 120÷90×60 + 5 ≈ 85 min
    Paris             & Crépy-en-Valois   & BUS    & 8   & 1{,}7  & 59  \\ % 80÷90×60 + 5 ≈ 59 min
    
    % Routes en AVION (800 km/h + 120 min temps'aéroport)
    Lille             & Paris             & AVION  & 120 & 85{,}0 & 136 \\ % 220÷800×60 + 120 ≈ 136 min
    Dunkerque         & Paris             & AVION  & 110 & 80{,}0 & 141 \\ % 280÷800×60 + 120 ≈ 141 min

    \bottomrule
  \end{tabular}
  \caption{Données complètes du réseau de transport}
\end{table}

{\small\itshape Les durées intègrent les temps d'attente habituels
(enregistrement aéroport pour l'avion, correspondances en gare).}

\newpage
Le scénario retenu pour ce prototype est le suivant :

\begin{itemize}
  \item Mode de transport : TRAIN uniquement
  \item Critère d'optimisation : le prix (on cherche à minimiser le coût en euros)
  \item Ville de départ : Lille
  \item Ville d'arrivée : Crépy-en-Valois
  \item Nombre d'itinéraires souhaités : 4
\end{itemize}

Ce scénario est représentatif d'un voyageur contraint par son budget, souhaitant
rejoindre Crépy-en-Valois depuis Lille en train. Le Tableau~2 présente les quatre
meilleurs itinéraires identifiés, classés par ordre croissant de prix.

\noindent\textbf{Tableau 2 : Les 4 meilleurs itinéraires TRAIN de Lille à
Crépy-en-Valois (critère : prix)}

\begin{table}[H]
  \centering
  \begin{tabularx}{\textwidth}{Xrr}
    \toprule
    \textbf{Itinéraire} & \textbf{Prix (€)} & \textbf{Classement} \\
    \midrule
    Lille → Cambrai → Paris → Crépy-en-Valois      & 61 & 1\textsuperscript{er} \\
    Lille → Dunkerque → Crépy-en-Valois            & 63 & 2\textsuperscript{e}  \\
    Lille → Paris → Crépy-en-Valois                & 67 & 3\textsuperscript{e}  \\
    Lille → Cambrai → Dunkerque → Crépy-en-Valois  & 81 & 4\textsuperscript{e}  \\
    \bottomrule
  \end{tabularx}
  \caption{Les 4 meilleurs itinéraires TRAIN de Lille à Crépy-en-Valois (critère : prix)}
\end{table}

{\small\itshape Les itinéraires empruntant l'avion ou le bus sont exclus du
scénario, le filtre portant exclusivement sur le mode TRAIN.}

% ──────────────────────────────
\section{Modèle pour l'exemple}
% ──────────────────────────────

Pour résoudre ce problème, on modélise le réseau filtré par un graphe orienté
valué. Chaque ville du réseau devient un sommet, et chaque connexion TRAIN devient
une arête orientée dont le poids correspond au critère choisi, ici le prix en euros.

\begin{figure}[H]
  \centering
  \includegraphics[width=0.60\textwidth]{imgrapport/graphesv1}
  \caption{Graphe orienté valué, scénario TRAIN, critère PRIX (Lille vers Crépy-en-Valois)}
\end{figure}

Dans ce graphe, les cinq sommets sont les villes
\{Lille, Dunkerque, Cambrai, Paris, Crépy-en-Valois\}. Les arêtes orientées et
valuées représentent les connexions TRAIN disponibles. À titre d'exemple,
l'arête Lille~→~Paris porte le poids~55 (prix en euros), l'arête
Lille~→~Cambrai porte le poids~14, et l'arête Paris~→~Crépy-en-Valois porte le
poids~12.

Les quatre meilleurs chemins dans ce graphe, exprimés sous forme de suites
d'arêtes avec leur poids cumulé, sont :

\begin{itemize}
  \item \textbf{Chemin~1 :} Lille → Cambrai → Paris → Crépy-en-Valois,
        poids total $= 14 + 35 + 12 = \mathbf{61}$~€
  \item \textbf{Chemin~2 :} Lille → Dunkerque → Crépy-en-Valois,
        poids total $= 18 + 45 = \mathbf{63}$~€
  \item \textbf{Chemin~3 :} Lille → Paris → Crépy-en-Valois,
        poids total $= 55 + 12 = \mathbf{67}$~€
  \item \textbf{Chemin~4 :} Lille → Cambrai → Dunkerque → Crépy-en-Valois,
        poids total $= 14 + 22 + 45 = \mathbf{81}$~€
\end{itemize}

L'algorithme KPCC explore les chemins sans cycle dans le graphe et les retourne
ordonnés par poids croissant. Les résultats ci-dessus ont été obtenus et vérifiés
par l'exécution de la classe \ic{GraphesV1} décrite en Section~1.4.

% ──────────────────────────────
\section{Modélisation pour la Version~1 dans le cas général}
% ──────────────────────────────

On se donne un problème de recherche d'itinéraire défini par : un ensemble de
villes~$V$, un ensemble de connexions~$C$ entre ces villes, une modalité de
transport~$m$ choisie, un critère d'optimisation~$c$ parmi
\{TEMPS, PRIX, CO2\}, une ville de départ~$s$, une ville d'arrivée~$t$, et un
entier~$k$ représentant le nombre d'itinéraires souhaités. Voici comment ce
problème se traduit en termes de graphes.

\noindent\textbf{Type de graphe}

On utilise un \emph{multigraphe orienté valué}. L'orientation est nécessaire car
une connexion de la ville~$x$ vers la ville~$y$ n'implique pas l'existence d'une
connexion en sens inverse : les lignes de transport sont directionnelles. La
valuation permet d'associer à chaque arête le coût correspondant au critère
choisi. La structure de multigraphe est justifiée par le fait que deux
villes pourraient théoriquement être reliées par plusieurs connexions de même
modalité, ce qui correspondrait à des arêtes parallèles dans le graphe.

\noindent\textbf{Sommets du graphe}

Les sommets du graphe sont les villes qui apparaissent dans au moins une
connexion de la modalité choisie, soit en départ, soit en arrivée. Autrement
dit, on ne retient que les villes réellement desservies par le mode de transport
sélectionné ; les autres sont ignorées. Dans notre exemple, filtrer sur TRAIN
exclut toute ville uniquement accessible en bus ou en avion.

\noindent\textbf{Arêtes du graphe et leurs poids}

Il existe une arête orientée du sommet~$x$ vers le sommet~$y$ dans le graphe
si et seulement si il existe une connexion de modalité~$m$ allant de~$x$ à~$y$
dans l'ensemble~$C$. Le poids de cette arête est la valeur du critère~$c$ pour
cette connexion :

\begin{itemize}
  \item si $c = \text{PRIX}$ → poids = prix en euros
  \item si $c = \text{TEMPS}$ → poids = durée en minutes
  \item si $c = \text{CO2}$ → poids = émissions en kg CO\textsubscript{2}e
\end{itemize}

Cette construction garantit que minimiser le poids total d'un chemin revient
exactement à minimiser le critère souhaité par l'utilisateur.

\noindent\textbf{Algorithme de résolution}

Une fois le graphe construit, on applique l'algorithme des $k$ plus courts chemins
(KPCC) en appelant \ic{AlgorithmeKPCC.kpcc(graphe,\,s,\,t,\,k)}. Cet algorithme
retourne une liste ordonnée d'au plus~$k$ objets \ic{Chemin}, chacun contenant la
liste des arêtes empruntées et le poids total du chemin. Les chemins sont sans
cycle et triés par poids croissant, ce qui correspond directement aux $k$ meilleurs
itinéraires selon le critère de l'utilisateur.

Cette modélisation est pensée pour permettre une évolution naturelle vers la
Version~2 : il suffira d'enrichir le graphe avec des nœuds de correspondance pour
prendre en charge les changements de modalité entre deux étapes d'un même voyage.

% ──────────────────────────────
\section{Prototype pour la Version~1}
% ──────────────────────────────

La classe \ic{GraphesV1} est une classe Java autonome et exécutable. Conformément
aux exigences du sujet, elle ne dépend d'aucune autre classe du projet, à
l'exception des interfaces \ic{Lieu} et \ic{Connexion} de la bibliothèque fournie.

La classe reproduit le scénario décrit en Section~1.1 : elle instancie les cinq
villes comme sommets, ajoute les connexions TRAIN avec leurs prix comme poids
d'arêtes, puis appelle \ic{AlgorithmeKPCC.kpcc} pour calculer les quatre meilleurs
itinéraires de Lille à Crépy-en-Valois. Les résultats sont affichés de façon
lisible, avec pour chaque itinéraire le rang, le coût total en euros et le détail
des étapes ville par ville.

\noindent\textbf{Lien vers le code source (commit avant la date limite) :}\\*
\href{https://gitlab.univ-lille.fr/sae2.01-2.02/2026/E5/-/blob/main/src/main/java/sae/transport/comparison/graphs/GraphesV1.java?ref_type=heads}{\ic{GraphesV1.java} sur GitLab (équipe E5)}

% ─────────────────────────────────────────────
\chapter{Version 2 : Plusieurs moyens de transport}
% ─────────────────────────────────────────────

% ──────────────────────────────
\section{Exemple}
% ──────────────────────────────

Pour illustrer la Version~2, nous réutilisons le même réseau de cinq villes du nord
de la France : Lille, Dunkerque, Cambrai, Paris et Crépy-en-Valois. Le réseau est
identique à celui de la Version~1 (Tableau~1). La nouveauté réside dans la possibilité
de combiner plusieurs modes de transport au sein d'un même itinéraire, moyennant un
coût de correspondance lorsqu'un changement de mode a lieu dans une ville.

\noindent\textbf{Tableau 1 : Données complètes du réseau de transport (inchangé par
rapport à la Version~1)}

\begin{table}[H]
  \centering
  \small
  \begin{tabular}{lllrrr}
    \toprule
    \textbf{Départ} & \textbf{Arrivée} & \textbf{Mode}
      & \textbf{Prix (€)} & \textbf{CO\textsubscript{2} (kg)}
      & \textbf{Durée (min)} \\
    \midrule
    Lille             & Dunkerque         & TRAIN  & 18  & 0{,}8  & 65  \\
    Dunkerque         & Lille             & TRAIN  & 18  & 0{,}8  & 65  \\
    Lille             & Cambrai           & TRAIN  & 14  & 0{,}7  & 55  \\
    Cambrai           & Lille             & TRAIN  & 14  & 0{,}7  & 55  \\
    Cambrai           & Paris             & TRAIN  & 35  & 1{,}8  & 95  \\
    Paris             & Cambrai           & TRAIN  & 35  & 1{,}8  & 95  \\
    Paris             & Crépy-en-Valois   & TRAIN  & 12  & 0{,}5  & 45  \\
    Crépy-en-Valois   & Paris             & TRAIN  & 12  & 0{,}5  & 45  \\
    Lille             & Paris             & TRAIN  & 55  & 2{,}4  & 60  \\
    Paris             & Lille             & TRAIN  & 55  & 2{,}4  & 60  \\
    Dunkerque         & Cambrai           & TRAIN  & 22  & 1{,}0  & 80  \\
    Cambrai           & Dunkerque         & TRAIN  & 22  & 1{,}0  & 80  \\
    Dunkerque         & Crépy-en-Valois   & TRAIN  & 45  & 2{,}0  & 150 \\
    Dunkerque         & Cambrai           & BUS    & 12  & 2{,}1  & 110 \\
    Cambrai           & Crépy-en-Valois   & BUS    & 18  & 2{,}8  & 130 \\
    Paris             & Crépy-en-Valois   & BUS    & 8   & 1{,}2  & 75  \\
    Lille             & Paris             & AVION  & 120 & 85{,}0 & 55  \\
    Dunkerque         & Paris             & AVION  & 110 & 80{,}0 & 50  \\
    \bottomrule
  \end{tabular}
  \caption{Données complètes du réseau de transport (inchangé par rapport à la Version~1)}
\end{table}

{\small\itshape Les connexions BUS et AVION, absentes du scénario V1, jouent ici un
rôle dans les itinéraires multi-modaux.}

Le Tableau~2 présente les coûts de correspondance retenus pour les villes où plusieurs
modes coexistent. Ces coûts reflètent les délais réels de changement de transport :
temps de déplacement entre les terminaux, enregistrement, attente en gare ou à
l'aéroport.

\noindent\textbf{Tableau 2 : Coûts de correspondance par ville et par paire de modes}

\begin{table}[H]
  \centering
  \begin{tabular}{lllrrr}
    \toprule
    \textbf{Ville} & \textbf{Mode 1} & \textbf{Mode 2}
      & \textbf{Durée (min)} & \textbf{CO\textsubscript{2} (kg)} & \textbf{Prix (€)} \\
    \midrule
    Lille  & TRAIN & AVION & 130 & 0{,}1 & 20 \\
    Lille  & TRAIN & BUS   & 20  & 0     & 0  \\
    Lille  & AVION & BUS   & 120 & 0{,}1 & 20 \\
    Paris  & TRAIN & AVION & 130 & 0{,}1 & 20 \\
    Paris  & TRAIN & BUS   & 15  & 0     & 0  \\
    Paris  & AVION & BUS   & 120 & 0{,}1 & 20 \\
    \bottomrule
  \end{tabular}
  \caption{Coûts de correspondance par ville et par paire de modes}
\end{table}

{\small\itshape Les correspondances sont symétriques : le coût de TRAIN~→~BUS est
identique à celui de BUS~→~TRAIN.}

Le scénario retenu pour ce prototype est le suivant :

\begin{itemize}
  \item Modes de transport : tous les modes disponibles (TRAIN, BUS, AVION)
  \item Critère d'optimisation : le prix (on cherche à minimiser le coût total en
        euros, correspondances incluses)
  \item Ville de départ : Lille
  \item Ville d'arrivée : Crépy-en-Valois
  \item Nombre d'itinéraires souhaités : 4
\end{itemize}

Le Tableau~3 présente les quatre meilleurs itinéraires identifiés, en tenant compte à
la fois des connexions de transport et des coûts de correspondance. On notera que les
deux meilleurs itinéraires sont multi-modaux, ce qui illustre bien l'apport de la
Version~2 par rapport à la Version~1.

\noindent\textbf{Tableau 3 : Les 4 meilleurs itinéraires multi-modaux de Lille à
Crépy-en-Valois (critère : prix)}

\begin{table}[H]
  \centering
  \small
  \begin{tabularx}{\textwidth}{Xrr}
    \toprule
    \textbf{Itinéraire} & \textbf{Prix (€)} & \textbf{Classement} \\
    \midrule
    Lille (TRAIN) → Cambrai (TRAIN) → Paris (TRAIN) [corresp.\ TRAIN→BUS]
      → Crépy-en-Valois (BUS) & 57 & 1\textsuperscript{er} \\
    Lille (TRAIN) → Cambrai (TRAIN) → Paris (TRAIN) → Crépy-en-Valois (TRAIN)
      & 61 & 2\textsuperscript{e} \\
    Lille (TRAIN) → Dunkerque (TRAIN) → Crépy-en-Valois (TRAIN)
      & 63 & 3\textsuperscript{e} \\
    Lille (TRAIN) → Paris (TRAIN) [corresp.\ TRAIN→BUS] → Crépy-en-Valois (BUS)
      & 63 & 4\textsuperscript{e} \\
    \bottomrule
  \end{tabularx}
  \caption{Les 4 meilleurs itinéraires multi-modaux de Lille à Crépy-en-Valois (critère : prix)}
\end{table}

{\small\itshape Les correspondances sont indiquées entre crochets. Le 1\textsuperscript{er}
itinéraire est multi-modal : TRAIN jusqu'à Paris, puis BUS jusqu'à Crépy.}

% ──────────────────────────────
\section{Modèle pour l'exemple}
% ──────────────────────────────

Pour prendre en charge la multi-modalité et les coûts de correspondance, on enrichit le
modèle de la Version~1 en dupliquant les sommets : chaque ville est représentée par
autant de sommets que de modes de transport disponibles dans cette ville. Ainsi, Lille
devient trois sommets distincts (\ic{Lille\_TRAIN}, \ic{Lille\_AVION} et
\ic{Lille\_BUS}), tandis que Dunkerque, qui n'est accessible qu'en train et en avion,
n'en possède que deux.

\begin{figure}[H]
  \centering
  \includegraphics[width=0.85\textwidth]{imgrapport/graphesv2}
  \caption{Graphe orienté valué pour le scénario V2 (multi-modal, critère PRIX)}
  \label{fig:graphesv2}
\end{figure}

Les arêtes de transport relient les sommets de même modalité entre deux villes
différentes : \ic{Lille\_TRAIN}~→~\ic{Paris\_TRAIN} avec le poids~55,
\ic{Paris\_TRAIN}~→~\ic{Crépy-en-Valois\_TRAIN} avec le poids~12,
\ic{Paris\_BUS}~→~\ic{Crépy-en-Valois\_BUS} avec le poids~8.

Les arêtes de correspondance relient les sommets de modes différents au sein d'une même
ville : \ic{Paris\_TRAIN}~→~\ic{Paris\_BUS} avec le poids~0 (aucun surcoût pour ce
changement à Paris). C'est précisément cette arête qui rend possible l'itinéraire~1
du Tableau~3, le plus économique.

Les quatre meilleurs chemins dans ce graphe, avec leurs poids cumulés, sont :

\begin{itemize}
  \item \textbf{Chemin~1 :} \ic{Lille\_TRAIN} → \ic{Cambrai\_TRAIN} → \ic{Paris\_TRAIN}
        → \ic{Paris\_BUS} → \ic{Crépy\_BUS},
        poids total $= 14 + 35 + 0 + 8 = \mathbf{57}$~€
  \item \textbf{Chemin~2 :} \ic{Lille\_TRAIN} → \ic{Cambrai\_TRAIN} → \ic{Paris\_TRAIN}
        → \ic{Crépy\_TRAIN},
        poids total $= 14 + 35 + 12 = \mathbf{61}$~€
  \item \textbf{Chemin~3 :} \ic{Lille\_TRAIN} → \ic{Dunkerque\_TRAIN} → \ic{Crépy\_TRAIN},
        poids total $= 18 + 45 = \mathbf{63}$~€
  \item \textbf{Chemin~4 :} \ic{Lille\_TRAIN} → \ic{Paris\_TRAIN} → \ic{Paris\_BUS}
        → \ic{Crépy\_BUS},
        poids total $= 55 + 0 + 8 = \mathbf{63}$~€
\end{itemize}

Ces résultats ont été obtenus et vérifiés par l'exécution de la classe \ic{GraphesV2}
décrite en Section~2.4.

% ──────────────────────────────
\section{Modélisation pour la Version~2 dans le cas général}
% ──────────────────────────────

L'idée centrale de la Version~2 est simple : puisqu'un voyageur peut changer de mode
de transport en cours de route, il faut que le graphe soit capable de représenter non
seulement les déplacements, mais aussi les changements de mode. Pour cela, on utilise
une technique appelée \textbf{duplication des sommets}.

\noindent\textbf{Type de graphe}

Comme pour la Version~1, on construit un \emph{multigraphe orienté valué}. La structure
est la même, seule la façon de créer les sommets et les arêtes change.

\noindent\textbf{Sommets du graphe}

Dans la Version~1, chaque ville était un sommet. Dans la Version~2, chaque ville devient
plusieurs sommets, un par mode de transport disponible dans cette ville. Par exemple,
Lille devient trois sommets : \ic{Lille\_TRAIN}, \ic{Lille\_AVION} et \ic{Lille\_BUS}.
Cambrai, accessible uniquement en train et en bus, ne donne que deux sommets :
\ic{Cambrai\_TRAIN} et \ic{Cambrai\_BUS}.

Cette distinction est indispensable : elle permet au graphe de savoir dans quel
«~état~» se trouve le voyageur à chaque étape, ce qui est nécessaire pour calculer le
coût d'un éventuel changement de mode.

\noindent\textbf{Arêtes de transport}

Une arête de transport relie deux sommets de même mode dans des villes différentes.
Par exemple, il existe une arête de \ic{Lille\_TRAIN} vers \ic{Paris\_TRAIN} avec le
poids~55, et une arête de \ic{Paris\_BUS} vers \ic{Crépy-en-Valois\_BUS} avec le
poids~8. Ces arêtes représentent les trajets classiques, exactement comme dans la
Version~1.

\noindent\textbf{Arêtes de correspondance}

Une arête de correspondance relie deux sommets de modes différents dans la même ville.
Par exemple, \ic{Paris\_TRAIN} vers \ic{Paris\_BUS} avec le poids~0. Ces arêtes
modélisent le changement de mode dans une ville ; sans elles, il serait impossible de
passer du train au bus en cours de route. Le sous-graphe formé par ces arêtes est un
sous-graphe induit par la relation de changement de mode.

\noindent\textbf{Algorithme de résolution}

On applique le même algorithme KPCC que pour la Version~1, en le lançant une fois par
mode disponible à la ville d'arrivée, puis en fusionnant les résultats pour garder les
$k$ meilleurs.

% ──────────────────────────────
\section{Prototype pour la Version~2}
% ──────────────────────────────

La classe \ic{GraphesV2} est une classe Java autonome et exécutable. Comme
\ic{GraphesV1}, elle ne dépend d'aucune autre classe du projet, à l'exception des
interfaces \ic{Lieu} et \ic{Connexion} de la bibliothèque fournie.

La classe met en œuvre la technique des sommets dupliqués décrite en Section~2.3 : les
cinq villes sont instanciées en plusieurs sommets selon leurs modes disponibles
(12~sommets au total), les arêtes de transport relient les sommets de même mode, et
les arêtes de correspondance relient les modes différents au sein de Lille et Paris
avec les coûts du Tableau~2.

L'algorithme KPCC est appelé deux fois, une fois vers \ic{Crépy-en-Valois\_TRAIN} et
une fois vers \ic{Crépy-en-Valois\_BUS}, et les résultats sont fusionnés puis triés
par poids croissant. L'affichage filtre les arêtes de correspondance pour ne montrer que
les points d'intérêt : départ, arrivée, et lieux de changement de modalité.

\noindent\textbf{Lien vers le code source (commit avant la date limite) :}\\*
\href{https://gitlab.univ-lille.fr/sae2.01-2.02/2026/E5/-/blob/main/src/main/java/sae/transport/comparison/graphs/GraphesV2.java?ref_type=heads}{\ic{GraphesV2.java} sur GitLab (équipe E5)}

\end{document}
